\chapter{Tips for Compiling Large Documents}\label{appx:compile-time}
\section{Compiling Chapters As Needed}
If you \verb|\include{}| your chapters, then you may speed up compilation considerably by using the \verb|\includeonly{}| command in your preamble.
Suppose you are working on some chapter which is sourced in the file \verb|chapters/file.tex|; then, putting the line \verb|\includeonly{chapters/file}| in the preamble of your document will tell \LaTeX\ to compile that chapter in isolation.
You are encouraged to read \href{https://tex.stackexchange.com/questions/246/when-should-i-use-input-vs-include}{this Stack Exchange post} about the use of \verb|\include|, \verb|\input|, and \verb|\includeonly|.

\section{Externalizing Graphics}
If you have a large number of \texttt{tikz} graphics, it is a good idea to \textit{externalize} these images so that they are not compiled every time you recompile your document.
There are a number of ways to do this, the most straightforward being to compile the figures somewhere else (\textit{e.g.}, in standalone \texttt{.tex} file found in your \verb|figs/| directory) and then to include the once-compiled graphics in your thesis using the \verb|\includegraphics| command.
If you are working on Overleaf, you are recommended to look into the TikZ library \href{https://tikz.dev/library-external}{external}, which handles this process automatically.
Overleaf itself has resources on how to externalize TikZ graphics, which can be found at \href{https://docs.overleaf.com/troubleshooting-and-support/fixing-and-preventing-compile-timeouts/reducing-the-compile-time-for-diagrams}{this link}.

\section{Other Miscellaneous Tips}
You are generally discouraged from modifying the \texttt{tamuthesis.cls} file, since doing so will likely result in a document which fails to meet Texas A\&M's formatting and accessibility requirements; however, an easy and seemingly harmless change for the curious is to replace the \texttt{unicode-math} package with \href{https://ctan.org/pkg/lua-unicode-math}{\texttt{lua-unicode-math}}.
The latter package is a modern replacement for \texttt{unicode-math} specific to Lua\TeX.
In our testing, it seems to work just as well and has the added benefit of making the document compile (marginally) faster.
In the future, this template might load \texttt{lua-unicode-math} by default; however, we presently defer to \texttt{unicode-math} as it is more stable and has been tested for longer.
